% Workshop paper draft. Target: AGENT 2027 (ICSE 2027 workshop on Agentic Engineering), full paper, up to 8 pages
% including references, IEEE conference template, single-anonymous. Deadline Nov 20, 2026 AoE.
% Numbers marked \pilot{} come from the 60-bug pilot and are replaced by fresh-bug results before submission.
\documentclass[10pt,conference]{IEEEtran}
\usepackage{booktabs}
\usepackage{xcolor}
\usepackage{cite}
\usepackage{url}

\newcommand{\todo}[1]{\textcolor{red}{[TODO: #1]}}
\newcommand{\pilot}[1]{\textcolor{blue}{#1}}

\begin{document}

\title{When Tests Lie: Cascades, Reviewers and Failure Attribution in Coding Agents}

\author{\IEEEauthorblockN{Nileet Savale}
\IEEEauthorblockA{Independent Researcher\\savalenileet@gmail.com}}

\maketitle

\begin{abstract}
\todo{Write last. One sentence each: setting (verify-and-retry agents gate on the repo's own tests; cascades send
work to a cheap model first and escalate on test failure); what we measure (how often test-passing fixes are wrong,
when test-gated cascades help, whether a reviewer is a better escalation trigger); data (SWE-bench Verified, $N$
fresh bugs, Claude and Gemini families); the three headline numbers; the takeaway for practitioners.}
\end{abstract}

\begin{IEEEkeywords}
LLM agents, program repair, model cascades, patch correctness, SWE-bench, cost
\end{IEEEkeywords}

\section{Introduction}
% ~1 page
Coding agents increasingly work in a verify-and-retry loop: they edit the code, run the repository's tests, and
retry until the tests pass~\cite{swebench}. Passing the tests is the loop's definition of done. The same signal is a
natural basis for saving money. In a \emph{cascade}, a cheap model attempts the task first, and the task escalates
to a strong model only when the cheap model's fix fails the tests~\cite{frugalgpt,kapoor2024}. If the cheap model
solves most tasks, the cascade should approach the strong model's success rate at a fraction of its cost.

That reasoning assumes the tests tell the truth. A gate can only escalate the failures it sees. A fix that passes the
repository's tests but does not resolve the issue is accepted and never escalated: the tests lie, and the cascade
believes them. This is test overfitting, long studied in automated program repair~\cite{smith2015,qi2015}, in a new
place: inside the escalation decision of a multi-model agent.

We measure how often this happens, what it costs, and what can be done about it, on SWE-bench
Verified~\cite{swebenchverified} with models from two families (Anthropic's Claude and Google's Gemini). SWE-bench's
hidden tests, which the agent never sees, serve as ground truth. We develop the agent and the reviewer on a
\pilot{60-bug} pilot and report the main results on the \todo{$N$} remaining bugs. Our contributions:
\begin{itemize}
\item \textbf{How often the tests lie (RQ1).} \pilot{28\%} of the fixes a cheap model's loop accepted fail the hidden
tests; we classify why, with two human annotators.
\item \textbf{When test-gated cascades help (RQ2).} A cheap model's failures are either \emph{visible}, ending with
failing tests, or \emph{silent}, accepted but wrong. A test gate acts only on visible failures. In our pilot, Claude
Haiku's failures were almost all silent and its cascade gained nothing; Gemini Flash-Lite's were mostly visible and
its cascade gained \pilot{7} bugs.
\item \textbf{A reviewer as the escalation trigger (RQ3).} A strong model reviewing the cheap model's test-passing
diff catches silent failures that the tests miss. We report catch rate, false-rejection rate and end-to-end cost, for
same-family and cross-family reviewers.
\item \textbf{Failure attribution (RQ4).} Some tests fail before the agent changes anything. Telling those apart from
failures the agent caused, and stopping when only they remain, cut the cost of such bugs \pilot{4.4-fold} in our pilot.
\item All runs, diffs, reviews and labels \todo{released at URL}.
\end{itemize}

\section{Background and Related Work}
% ~1 page. Every citation to be checked against its source before submission.
\textbf{Plausible but incorrect patches.} Automated program repair has long distinguished \emph{plausible} patches,
which pass the tests used during repair, from correct ones. Smith et al.\ showed that patches fitted to one test suite
often break held-out tests~\cite{smith2015}, and Qi et al.\ found that most plausible patches from early repair tools
were wrong~\cite{qi2015}. For LLM agents on SWE-bench, Ahmed and Ganhotra measure overfitting to reproduction tests
the agent generates, at 21.8--33\%, and find that iterative refinement raises it~\cite{ahmed2025overfitting}. Other
work strengthens SWE-bench's own tests: PatchDiff finds that many ``solved'' patches behave differently from the
reference~\cite{wang2025patchdiff}, and UTBoost adds tests to instances whose tests are
insufficient~\cite{utboost}. Our gate is weaker still: the repository's existing tests, which is what a deployed
agent actually has.

\textbf{Cascades and routing.} FrugalGPT cascades LLMs with a learned scorer~\cite{frugalgpt}; Dekoninck et al.\ show
that the quality estimator decides whether cascading pays~\cite{dekoninck2025}. For code, Kapoor et al.\ ran an
escalation baseline that moved to a stronger model when example tests failed, and noted that cheap models' solutions
can pass the example tests and fail the evaluation tests~\cite{kapoor2024}. Chen et al.\ cascade code models on
agreement with generated tests~\cite{chen2024cascade}. Sun et al.\ find that escalating reasoning effort on failure
pays when the gate is the benchmark's own scoring tests~\cite{sun2026effort}. We study the deployable case: the gate
is the repository's tests, and the question is which failures it can see.

\textbf{LLMs as patch judges.} Closest to our reviewer, Cambronero et al.\ gate an agent's patches at Google with an
LLM validator, rejecting \todo{58\%} of bad patches at a cost of rejecting about a quarter of good
ones~\cite{cambronero2025}. Re-evaluations of patch-overfitting detectors find that many do no better than random
selection~\cite{williams2026}, so we compare our reviewer with a random rejector. To our knowledge, no prior work uses
a strong model's review of a cheap model's \emph{test-passing} patch as the trigger for escalation, measures the
end-to-end cost of doing so, or compares reviewers from the same and a different model family.

\textbf{Failure attribution.} Benchmarks usually remove instances whose tests fail even with the reference
fix~\cite{tddbench}. A deployed agent cannot: it must decide whether a failing test is its fault. We measure what that
decision saves inside the loop.

\section{Study Design}
% ~1.5 pages

\subsection{Research Questions}
\begin{description}
\item[RQ1] How often does a cheap model's fix that passes the repository's tests fail SWE-bench's hidden tests, and why?
\item[RQ2] When does a cascade that escalates on test failure improve on the cheap model alone, and at what cost?
\item[RQ3] Is a reviewer model a better escalation trigger than the tests? How do same-family and cross-family
reviewers differ?
\item[RQ4] Does telling failures the agent caused from failures it did not (tests already failing before any change)
cut wasted retries, without losing fixes?
\end{description}

\subsection{Benchmark and Ground Truth}
We use SWE-bench Verified~\cite{swebenchverified}: 500 human-validated Python issues from 12 repositories, each with
a gold patch and hidden tests (FAIL\_TO\_PASS and PASS\_TO\_PASS). The agent never sees the hidden tests; it sees the
issue text and the repository, and may run the repository's existing tests. A run is \emph{resolved} if its final
diff passes the hidden tests in the official SWE-bench harness. \todo{Pilot: 60 bugs, how chosen. Main set: the
$N$ remaining bugs, never used while designing the loop or the reviewer prompt.}

Before any agent runs, we run each bug's existing tests on the original code and with the gold patch applied, to
record which tests already fail and which fail even with the correct fix. The loop uses this to tell failures it
caused from failures it did not (Section~\ref{sec:loop}). \todo{Say how many bugs have tests that fail with the gold
patch.}

\subsection{The Verify-and-Retry Loop}
\label{sec:loop}
One coder model works in a checkout of the repository with tools to search, read and edit files. After each attempt
the loop runs the repository's tests that relate to the change and hands the failures back. An attempt has at most 25
tool steps and a run at most 5 attempts. The loop stops when the tests pass, when the only failing tests were already
failing before any change, after two attempts in a row with no changes, or when the same failure repeats.
\todo{Name the remaining lean settings that matter: retry context continues the conversation; short test output; a
low-effort test reviewer for proposed test changes; file listing capped at 100 entries; keyword pointers to likely
files. Per-run budget and time limit (10 min Claude, 20 min Gemini).}
As a sanity check that the loop is a realistic agent rather than a weak stand-in, on the pilot it resolved
\pilot{50 of 60} bugs with Claude Sonnet 5.5, against \pilot{45} for Claude Code with the same model (exact McNemar
$p=0.18$). On ImpossibleBench's ``conflicting'' tasks~\cite{impossiblebench}, whose visible tests contradict the issue,
it changed a test to pass on \pilot{0 of 8}.

\subsection{Setups}
\begin{itemize}
\item \textbf{Cheap alone}: the loop with the cheap model.
\item \textbf{Strong alone}: the loop with the strong model.
\item \textbf{Test-gated cascade}: the cheap model first writes notes on where the bug is, then runs the loop. If the
run does not end with passing tests, its changes are undone and the strong model runs the loop with the notes and a
summary of what was tried.
\item \textbf{Reviewed cascade}: as the test-gated cascade, but a cheap fix that passes the tests is also shown to a
reviewer (the issue, the diff, and the original code 40 lines around each changed hunk). If the reviewer rejects it,
the task escalates with the reviewer's reason. We wrote two reviewer prompts: v1, a plain accept-or-escalate prompt
used throughout the pilot, and v2, which also asks whether the fix missed a sibling code path or bent an existing test
to match itself. v2 was written after reading the pilot's failures, so it is evaluated only on the main set, where it
was fixed before any run; v1 is reported only on the pilot.
\end{itemize}

\subsection{Models}
\begin{tabular}{@{}lll@{}}
\toprule
Family & Cheap & Strong \\
\midrule
Anthropic & Claude Haiku 5.5 & Claude Sonnet 5.5 \\
Google & Gemini 3.5 Flash-Lite & Gemini 3.8 Flash \\
\bottomrule
\end{tabular}

\noindent Reviewers: Claude Sonnet 5.5 and Gemini 4 Carbon, each reviewing both families' cheap fixes. All models at
their default reasoning effort. \todo{Why Gemini 3.8 Flash rather than Carbon as the strong coder (cost).}

\subsection{Measures}
Resolved rate with 95\% Wilson intervals~\cite{bowyer2025}; paired comparisons per bug (McNemar).
Cost in US dollars from token counts at list prices, including cache reads. For RQ1, the false-accept rate is the share
of runs the loop accepted whose fix fails the hidden tests. For RQ3, catch rate (rejected bad fixes / bad fixes),
false-rejection rate (rejected good fixes / good fixes), MCC, and a random rejector with the same rejection rate as a
baseline~\cite{williams2026}. Failure causes for RQ1 were labelled by two annotators \todo{who; agreement $\kappa$}.

\section{Results}
% ~2.5 pages

\subsection{RQ1: Test-Passing but Wrong}
\pilot{Pilot: of 165 fixes Haiku's loop accepted, 46 (28\%) fail the hidden tests; Flash-Lite: 18 of 53 (34\%).}
\todo{Fresh-bug numbers; table of failure causes (pilot, 44 cases: unstated detail in hidden test 19, incomplete /
missed sibling 16, wrong location or symptom patch 4, evaluation artefact 5); compare with published overfitting rates.}

\subsection{RQ2: When Test-Gated Cascades Help}
\pilot{Pilot. Claude, 59 bugs: Haiku 43, cascade 42, Sonnet 50. Gemini, 41 bugs: Flash-Lite 17 (\$0.06/bug),
cascade 24 (\$0.29), 3.8 Flash 30 (\$0.56).} The cascade helps only when the cheap model's failures are visible to the
tests. We state this as an observation in two model families, not a general law. \todo{Fresh-bug table;
accuracy--cost Pareto plot.}

The unifying measure is each cheap model's split of failures into \emph{visible} (the loop ends with failing
tests, so a test gate escalates) and \emph{silent} (the loop accepts a fix that fails the hidden tests, so only a
reviewer can escalate). \todo{Figure: per family, resolved / visible failure / silent failure as stacked bars, cheap
model alone.} \pilot{Pilot, cheap model alone: of Haiku's 16 failures on 59 bugs, 15 are silent and 1 visible; of
Flash-Lite's 24 failures on 42 bugs, 10 are silent and 14 visible. A test gate can only act on the visible ones, which
is why it helps Flash-Lite and not Haiku.}

\subsection{RQ3: A Reviewer as the Escalation Trigger}
\pilot{Pilot, prompt v1. Sonnet on Haiku's fixes: catches 25/46, wrongly rejects 33/119. Carbon on Haiku's: 41/46,
58/119. Sonnet on Flash-Lite's: 11/18, 4/35. Carbon on Flash-Lite's: 17/18, 18/35.} \todo{Reviewed cascade, live
on the main set: resolved and cost against strong alone; random-rejector baseline; same- vs cross-family.}

\subsection{RQ4: Failure Attribution}
% Keep to about a third of a column.
\pilot{Pilot, 60 bugs, Sonnet 5.5. With attribution the loop resolved 50 at \$0.061/bug; with every failing test
counted as the agent's, 52 at \$0.114 (difference in resolved within noise, $p=0.63$). On the 14 bugs whose tests
already fail before any change, attribution cost \$0.83 against \$3.67 (1.0 vs 2.8 attempts). The stop on
``only pre-existing failures left'' fired 12 times; 11 of those runs resolved the bug.} \todo{Repeat on the main set;
how many of the $N$ bugs have pre-existing failures; the one loss (django-13837, stopped to flag a test for a human).}

\section{Discussion}
% ~0.5 page
\todo{Per-token price misleads: Gemini 3.8 Flash costs about 12x Sonnet per bug despite a lower token price (more
steps, more thinking, weaker implicit caching). Failure attribution: stopping when only pre-existing failures remain.
What a practitioner should do. Effort ladders: escalating from low to
default reasoning effort on test failure resolved \pilot{43} at \$0.067/bug vs \pilot{50} at \$0.062 for default
effort alone, the same blind spot as RQ2 on another axis; Sun et al.\ found ladders pay when the gate is the
benchmark's own scoring tests.}

\section{Threats to Validity}
% ~0.5 page
\todo{One benchmark, Python only, possibly contaminated (Verified); two model families, no open-weight model; single
runs (repeats on a subset); the pilot was run on emulated x86 images on Apple Silicon, the main set natively; the v2
reviewer prompt was shaped on the pilot, so only the main set tests it; the strong Gemini model is not the strongest
available; cost at list prices of November 2026.}

\section{Conclusion}
\todo{Two or three sentences.}

\section*{Data Availability}
\todo{Runs, diffs, reviews and scripts: repository link once the user decides to release.}

\section*{Acknowledgment}
\todo{Disclose AI assistance in drafting text, as IEEE requires.}

\bibliographystyle{IEEEtran}
\bibliography{refs}

\end{document}
